\documentclass{article}
\usepackage{tikz}
\usetikzlibrary{arrows.meta}
\usetikzlibrary{shapes.geometric}
\usetikzlibrary{positioning}
\usetikzlibrary{calc}

% Layer declarations -- required before \input-ing any TikZiT-produced .tikz file.
% "main" must be listed explicitly or content drawn outside nodelayer/edgelayer
% silently fails to render -- see book/chapter08_advanced_tikz.md.
\pgfdeclarelayer{nodelayer}
\pgfdeclarelayer{edgelayer}
\pgfsetlayers{main,nodelayer,edgelayer}

\input{../styles/epiflow.tikzstyles}

\begin{document}

\title{Example Manuscript with Embedded TikZiT Figures}
\author{}
\date{}
\maketitle

\section{Methods}

Participant flow through the study is shown in Figure~\ref{fig:flow}.
Eligibility screening with a labeled decision branch is shown in
Figure~\ref{fig:decision}. A curved-edge example illustrating the cookbook
technique is shown in Figure~\ref{fig:curved}.

\begin{figure}[htbp]
    \centering
    \begin{tikzpicture}
    \begin{pgfonlayer}{nodelayer}
        \node [style=include] (0) at (0, 0) {Assessed for eligibility (n=120)};
        \node [style=exclude] (1) at (4, -1.5) {Excluded (n=15)};
        \node [style=include] (2) at (0, -3) {Enrolled (n=105)};
    \end{pgfonlayer}
    \begin{pgfonlayer}{edgelayer}
        \draw [style=arrow] (0) to (1);
        \draw [style=arrow] (0) to (2);
    \end{pgfonlayer}
\end{tikzpicture}

    \caption{Participant flow.}
    \label{fig:flow}
\end{figure}

\begin{figure}[htbp]
    \centering
    \begin{tikzpicture}
    \begin{pgfonlayer}{nodelayer}
        \node [style=include] (0) at (0, 0) {Screened (n=200)};
        \node [style=decision] (1) at (0, -2.5) {Meets inclusion criteria?};
        \node [style=include] (2) at (4, -5) {Enrolled (n=150)};
        \node [style=exclude] (3) at (-4, -5) {Excluded (n=50)};
    \end{pgfonlayer}
    \begin{pgfonlayer}{edgelayer}
        \draw [style=arrow] (0) to (1);
        \draw [style=arrow] (1) to node [style=edgelabel] {Yes} (2);
        \draw [style=arrow] (1) to node [style=edgelabel] {No} (3);
    \end{pgfonlayer}
\end{tikzpicture}

    \caption{Eligibility screening with a labeled decision branch.}
    \label{fig:decision}
\end{figure}

\begin{figure}[htbp]
    \centering
    \begin{tikzpicture}
    \begin{pgfonlayer}{nodelayer}
        \node [style=process] (0) at (0, 0) {Assessed for eligibility (n=300)};
        \node [style=exclude] (1) at (4, -1.5) {Excluded (n=50)};
        \node [style=process] (2) at (0, -5) {Enrolled (n=250)};
        \node [style=process] (3) at (-4, -1.5) {Re-screened later (n=10)};
    \end{pgfonlayer}
    \begin{pgfonlayer}{edgelayer}
        \draw [style=arrow] (0) to (1);
        \draw [style=arrow] (0) to (2);
        \draw [style=arrow, bend left=20] (1) to node [style=edgelabel, sloped] {if re-contacted} (3);
        \draw [style=arrow, bend right=20] (3) to (0);
    \end{pgfonlayer}
\end{tikzpicture}

    \caption{Curved-edge example: a node with two outgoing paths using
             \texttt{bend left}/\texttt{bend right}.}
    \label{fig:curved}
\end{figure}

\end{document}
